\section{Project Goal}
SimpleSFT estimates the memory requirements for supervised fine-tuning (SFT) jobs
on popular large language models (LLMs) used in academic research. These models
include open-weight Qwen-family models, OLMo, LLaMA, and Mistral, which are common in current research
workflows and span sizes where a small change in distributed strategy can
determine whether a job runs at all.

\begin{quote}
Given a model, context length, resource pool, and training strategy, estimate
the expected per-rank peak memory and identify the training phase that produces
that peak.
\end{quote}

Today, users face slow trial-and-error development workflows to find parallel
strategies that meet their specific needs: short debug allocations on OSC
allow only one GPU, while larger allocations can require long queues or long reservations.
SimpleSFT helps users choose a plausible strategy before occupying GPUs for extended setup runs.

The state of OSC queues is often unpredictable and frustrating. This has led to a culture of
overprovisioning, where users request more resources than they expect to need just to be
sure they will have time to debug their jobs. This is inefficient for everyone: users wait
longer in queues, and the cluster is less available for other work. SimpleSFT provides a
more informed way to request resources, which can reduce wait times and improve overall
cluster efficiency.

This framing matters on shared academic clusters. An NLP researcher seeking to do SFT often begins
with a concrete scientific question, a model family, and some training data. The act of
engineering a system to run SFT efficiently is incidental to the true research question.
SimpleSFT lightens that work by providing a framework for getting to model training with less trial and error.

\section{System Overview}
The project has four main pieces:
\begin{itemize}[leftmargin=1.4em]
\item \textbf{Model inspection:} Hugging Face models are summarized as a
\texttt{ModelSpec}: layer count, hidden and intermediate sizes, vocabulary,
attention layout, parameter shapes, trainable linear layers, tensor-layout
roles, and optional architecture metadata.
\item \textbf{Analytical estimation:} an \texttt{EstimatorConfig} describes the
training strategy. The estimator builds resident state, retained activations,
workspace windows, and phase peaks, then returns a \texttt{MemoryResult}.
\item \textbf{Measurement and saved comparisons:} CUDA microsteps provide
ground-truth memory measurements for selected configurations. The saved
measurements let later estimator changes be checked against the same cases.
\item \textbf{Workflow tools:} CLI and local web paths expose estimate, search,
measure, compare, benchmark, report, YAML export, and SFT launch commands.
\end{itemize}
